\chapter{Rudiments de théorie des modèles}

\section{Structures et théories}

\subsection{Langages et structures}

\begin{defi}[Langage]
Un \kw{langage} $\cL$ est la donnée de:
    \begin{enumerate}[label= (\roman*)]
        \item Un ensemble $\cF$ de \kw{fonctions} et une arité $n_f \in \bNe$ pour chaque $f \in \cF$.
        \item Un ensemble $\cR$ de \kw{relations} et une arité $n_R \in \bNe$ pour
        chaque $R \in \cR$.
        \item Un ensemble $\cC$ de \kw{constantes}.
    \end{enumerate}
Les éléments de ces ensembles sont appelés les \kw{symboles} de $\cL$.
Rappelons que l'arité d'un symbole est le nombre d'arguments qu'il prend.
Par convention, les constantes sont d'arité 0.
\end{defi}

\begin{defi}[Structure]
Une \kw{$\cL$-structure} $\cM$ est la donnée d'un ensemble non vide $M$ appelé
\kw{domaine} et des \kw{interprétations} des symboles de $\cL$:
    \begin{enumerate}[label= (\roman*)]
        \item Une fonction $f^{\cM}:M^{n_f} \mapsto M$ pour chaque $f\in \cF$.
        \item Une ensemble $R^{\cM} \subset M^{n_R}$ pour chaque $R \in \cR$.
        \item Un élément $c^{\cM} \in M$ pour chaque $c \in \cC$.
    \end{enumerate}
\end{defi}

\begin{nota}
Dans la suite on notera, sauf mention contraire, les structures avec une lettre
majuscule cursive ($\cM,\cN,\cA, \ldots$) et leur domaine avec la même lettre
majuscule droite ($M,N,A, \ldots$).
\end{nota}

\begin{exem}[Langage des groupes]
L'étude des groupes induit le langage $\cL_{gp} =
(\cdot,{}^{-1},e)$ où $\cdot$ est une fonction binaire ${}^{-1}$ une fonction unaire et $e$ une constante.
Une $\cL_{gp}$-structure $\cG = (G,\cdot^{\cG},{}^{{-1}^{\cG}},e^{\cG})$ sera un ensemble $G$
muni d'une opération binaire $\cdot^{\cG}$, d'une opération unaire
${}^{{-1}^{\cG}}$ et d'un élément distingué $e^{\cG}$. \\
Par exemple $(\bR,\cdot,{}^{-1},1)$ est une $\cL_{gp}$-structure où on interprète
$\cdot$ comme la multiplication usuelle, ${}^{-1}$ comme l'inversion usuelle et
$e$ comme $1$.  \\
Remarquons qu'il n'y a pour le moment aucune règle
sémantique à vérifier, par exemple $\cN = (\bN,+,s,0)$ avec $\cdot^{\cN} = +$ l'addition usuelle, ${}^{{-1}^{\cN}}
= s$ le successeur et $e^{\cG} = 0$ est bien une
$\cL_{gp}$-structure même si ce n'est évidemment par un groupe.
\end{exem}

\begin{defi}[Sous-langage, expansion]
Soient $\cL,\cL'$ deux langages. Si $\cL \subset \cL'$, on dit que $\cL$ est
un \kw{sous-langage} de $\cL'$ et $\cL'$ une \kw{expansion} de $\cL$.
Soit $\cM'$ une $\mathcal{L'}$-structure, en restreignant l’interprétation donnée par $\mathcal{M'}$ au langage $\mathcal{L}$,
on obtient une $\mathcal{L}$-structure notée $\mathcal{M}'\upharpoonright_{\mathcal{L}}$. On dit que $\mathcal{M}'\upharpoonright_{\mathcal{L}}$
est une réduction de $\mathcal{M}'$, et $\mathcal{M}'$ une expansion de $\mathcal{M}'\upharpoonright_{\mathcal{L}}$.
\end{defi}

\begin{exem}[Langage des anneaux]
Considérons le langage des anneaux $\cL_{ring} =
(+,\cdot, - , 0 ,1 ) $, alors le langage $\cL_{agp} = (+,-,0 )$ des groupes
notés additivement est un \kw{sous-langage} de $\cL_{ring}$ car chaque symbole de
$\cL_{agp}$ et aussi un symbole de $\cL_{ring}$ avec la même arité. On dit aussi
que $\cL_{ring}$ est une \kw{expansion} de $\cL_{agp}$. En considérant la structure
$\bR_{ring} = (\bR;+,\cdot, - , 0 ,1 )$, alors $\bR_{agp} = (\bR;+,-,0)$ est
une réduction de $\bR_{ring}$. Si on change le domaine, on parle de
sous-structure et d'extensions, ainsi $\bZ_{ring} = (\bZ;+,\cdot, - , 0 ,1 )$
est une \kw{sous-structure} de $\bR_{ring}$ et $\bR_{ring}$ une \kw{extension} de
$\bZ_{ring}$.
\end{exem}

\begin{defi}[Plongement, sous-structure, extension]\label{plongement}
    Soient $\cM, \cN$ des $\cL$-structures, un $\cL$-plongement $\eta : \cM \rightarrow \cN$ (ou simplement \kw{plongement})
    est une fonction injective $\eta : M \rightarrow N$ qui préserve les
    interprétations de tous les symboles de $\cL$, plus précisément :
    \begin{enumerate}[label= (\roman*)]
        \item Pour chaque fonction $f$ de $\cL$ et tout $a_1, \ldots, a_{n_f} \in
        M$, \[\eta(f^{\cM}(a_1,\ldots,a_{n_f})) =
        f^{\cN}(\eta(a_1),\ldots,\eta(a_{n_f})) \]
        \item Pour chaque relation $R$ de $\cL$ et tout $a_1, \ldots, a_{n_R}
        \in M$, \[(a_1, \ldots, a_{n_R}) \in R^{\cM} \iff (\eta(a_1), \ldots,
        \eta(a_{n_R})) \in R^{\cN} \]
        \item Pour chaque constante $c$ de $\cL$, $ \eta(c^{\cM}) = c^{\cN}$
    \end{enumerate} 

    Si un plongement est bijectif, on parle d'\kw{isomorphisme}. On a alors un
    plongement $\mu : \cN \rightarrow \cM$ tel que $\eta \circ \mu$ est
    l'identité de $\cN$ et $\mu \circ \eta$
    l'identité de $\cM$, on le note, comme d'habitude, $\eta^{-1}$.
    Si de plus, $\cN = \cM$, on parle d'\kw{automorphisme}.

    Si l'injection canonique est un plongement, $\cM$ est une \kw{sous-structure} de $\cN$ et $\cN$ une \kw{extension} de $\cM$.
\end{defi}

\begin{exem}
    \begin{itemize}
        \item Les injections évidentes $\bZ \inj{} \bQ \inj{} \bR \inj{} \bC$ fournissent les
$\cL_{ring}$-plongements \[ \bZ_{ring} \inj{} \bQ_{ring} \inj{} \bR_{ring}
\inj{} \bC_{ring} .\] 
        \item Si $\eta : \bZ \rightarrow \bR$ est la fonction $\eta(x) = e^x$, alors $\eta$
est un $\cL_{gp}$-plongement de $(\bZ,+,0)$ dans $(\bR,\cdot,1)$.
        \item Les seuls automorphismes de $\bZ_{agp}$ sont l'identité et la
        fonction $x \mapsto -x$. En effet, si $\eta : \bZ \rightarrow \bZ$ est
        un plongement, alors $\eta(0) = 0$ et si $\eta(1) = n$, pour tout $m \in
        \bNe$, on a \[ \eta(m)=\eta(\underbrace{1+\cdots+1}_\text{m}) =
        \underbrace{\eta(1)+\cdots+\eta(1)}_\text{m} = mn .\] et de même,
        $\eta(-n)=-\eta(n)=-mn$, pour avoir un inverse, $\eta$ doit être
        surjective et donc on doit avoir $n = \pm 1$.
    \end{itemize}

\end{exem}


\subsection{Formules et modèles}

En plus des symboles spécifiques à un langage définis plus haut (relations,
fonctions et constantes), nous travaillons avec des symboles communs à tout langage:
les symboles logiques $=,\wedge, \neg, \exists, (,) $, et les variables
($x,y,z,\ldots$). Nous donnons le formalisme nécessaire pour construire des chaînes de symboles 
via un langage donné (appellées termes et formules) et leurs interpretations dans une structure.

\begin{defi}[Terme]\label{defi-terme}
On définit l'ensemble des termes d'un langage $\cL$ de manière récursive:
\begin{enumerate}[label= (\roman*)]
    \item Chaque variable est un terme.
    \item Chaque constante de $\cL$ est un terme.
    \item Si $f$ est une fonction de $\cL$ d'arité $n$, et $t_1, \ldots, t_n$
    des termes, alors $f(t_1,\ldots, t_n)$ est un terme.
    \item Tout expression obtenue exclusivement par l'application d'un nombre
    fini de fois des règles précédentes est un terme.
\end{enumerate}

Si $\cM$ est une $\cL$-structure, $t$ un terme et $\overline{x} = (x_1, \ldots,
x_n)$ des variables incluant toutes celles qui apparaissent dans $t$. La
définition récursive des termes permet la définition récursive d'une fonction
$t^{\cM} : M^n \rightarrow M$ comme suit:
\begin{enumerate}[label= (\roman*)]
    \item Si $t=c$ avec $c$ une constante, $t^{\cM}(\overline{x}) = c^{\cM}$.
    \item Si $t=x_i$, alors $t^{\cM}(\overline{x}) = x_i$.
    \item Si $t=f(t_1, \ldots, t_r)$, alors $t^{\cM}(\overline{x}) =
    f^{\cM}(t_1^{\cM}(\overline{x}),\ldots, t_r^{\cM}(\overline{x}))$
\end{enumerate}
\end{defi}

\begin{defi}[Formule]\label{termeformule}
On définit l'ensemble des \kw{formules} de $\cL$ de manière récursive:
\begin{enumerate}[label= (\roman*)]
    \item Si $t_1$ et $t_2$ sont des termes, alors $(t_1 = t_2)$ est une formule.
    \item Si $t_1, \ldots, t_n$ sont des termes et $R$ est une relation d'arité
    $n$, alors $R(t_1,\ldots,t_n)$ est une formule.
    \item Si $\varphi$ et $\psi$ sont des formules, alors $\neg \phi$ et
    $(\varphi \vee \psi)$ sont des formules.
    \item Si $\varphi$ est une formule et $x$ est une variable, alors $\exists x [\varphi]$ est une formule.
    \item Tout expression obtenue exclusivement par l'application d'un nombre
    fini de fois des règles précédentes est une formule.
\end{enumerate}
Les formules des points $(i)$ et $(ii)$ sont dites \kw{atomiques} car elles
ne contiennent pas de formules plus petites.
\end{defi}

\begin{exem}[Termes, formules]
Reprenons le langage des anneaux $\cL_{ring}$, les chaines \[ 0;
(1+1); -((1+1)+1); x\cdot(1+1); (x+y) \] sont des \kw{termes}.
On peut les interpréter dans $\bR_{ring}$ respectivement comme les nombres
$0;2;-3$ et les fonctions $x \mapsto 2x$ et $(x,y) \mapsto x+y$.
Notons qu'un terme ne dépend que du langage mais son interprétation dépend de la structure.
Les \kw{formules} sont des chaines qui donnent des informations sur la structure ou
ses éléments via une valeur de vérité qu'on leur associe une fois interpretées.
Par exemple $(1+1) = 0$ est une \kw{formule} qui  est \kw{fausse} dans $\bR_{ring}$
mais \kw{vraie} dans $\bZnZ{2}_{ring}$, $\exists x [(x\cdot x) = (1+1)]$ affirme que
$2$ est un carré, elle est fausse dans $\bZ_{ring}$ mais vraie dans $\bR_{ring}$.
\end{exem}

\begin{defi}[Variable libre/liée, phrase]
Une variable $x$ d'une formule est \kw{quantifiée} si elle suit immédiatement le
symbole $\exists$, on dit alors que le quantificateur $\exists x$ \kw{lie} la
variable. Une variable qui n'est pas liée est dite \kw{libre}.
Une formule sans variable libre est appelée une \kw{phrase} ou un \kw{énoncé}.
\end{defi}

\begin{exem}
Dans la formule $x < 1 + 1$, la variable $x$ est libre: la formule donne une information directe sur la variable $x$ considérée. Tandis que 
dans la formule $ \exists x [x \cdot x = 2]$, la variable $x$ est quantifiée: l'information donnée est sur la structure considérée plutot que sur une certaine variable, 
en effet, cette formule équivaut $ \exists y [y \cdot y = 2]$.
\end{exem}

\begin{rema}[Abréviations]
Notre langage formel est très restreint, nous utiliserons des abréviations usuelles
pour alléger les formules. Nous écrirons:

\begin{alignat*}{2}
\text{différent: } &(t_1 \neq t_2) && \text{ pour } \neg(t_1 = t_2), \\
\text{ou: } &(\varphi \vee \psi) && \text{ pour } \neg(\neg\varphi \wedge \neg\psi), \\
\text{implication: } &(\varphi \rightarrow \psi) && \text{ pour } (\neg\varphi \vee \psi), \\
\text{équivalence: } &(\varphi \leftrightarrow \psi) && \text{ pour } ((\varphi \rightarrow \psi) \wedge (\psi \rightarrow \varphi)), \\
\text{pour tout: } &\forall x[\varphi] && \text{ pour } \neg{\exists}x[\neg\varphi].
\end{alignat*}

On préfère restreindre notre langage au strict minimum tout en utilisant ces
autres symboles comme de simples abréviations pour ne pas alourdir nos preuves.
En effet, comme le quantificateur universel ou la disjonction sont définis comme 
des combinaisons de symboles de notre langage restreint, nous pourrons les utiliser
dans nos formules sans avoir à traiter leurs cas spécifiques lors des définitions ou des preuves par induction. 
Ces abréviations sont exploitées dès la définition suivante.
\end{rema}

\begin{defi}
Soient $\cM$ une $\cL$-structure, $\phi$ une $\cL$-formule prenant ses variables libres dans $ \vbar = (v_1,\dots,v_m)$
et soit $\abar \in M^n$. On définit de manière récursive $\cM \models \phi(\abar)$:
\begin{enumerate}
\item Si $\phi$ est $t_1 = t_2$, alors $\cM \models \phi(\abar)$ si $t_1^{\cM}(\abar) = t_2^{\cM}(\abar)$
\item Si $\phi$ est $\cR (t_1,\dots,t_{n_{\cR}})$, alors $\cM \models \phi(\abar)$ si $(t_1^{\cM}(\abar),\dots,t_{n_{\cR}}^{\cM}(\abar)) \in \cR^{\cM}$
\item Si $\phi$ est $\neg \psi$, alors $\cM \models \phi(\abar)$ si $\cM \not\models \psi(\abar)$
\item Si $\phi$ est $(\psi_1 \wedge \psi_2)$, alors $\cM \models \phi(\abar)$ si $\cM \models \psi_1(\abar)$ et $\cM \models \psi_2(\abar)$
\item Si $\phi$ est $\exists w \psi(\vbar,w)$, alors $\cM \models \phi(\abar)$ s'il existe $b \in M$ tel que $\cM \models \psi(\abar,b)$
\end{enumerate}
Si $\cM \models \phi(\abar)$, on dit que $\cM$ \kw{satisfait} $\phi(\abar)$, ou encore que $\phi(\abar)$ est \kw{vraie} dans $\cM$.
\end{defi}

\noindent La démonstration suivante donne une première mise en oeuvre du raisonnement par induction.

\begin{prop}
Soient $\cM, \cN$ deux $\cL$-structures telles que $\cM$ est une sous-structure de $\cN$. Soient $\abar \in M$ et
$\phi$ une $\cL$-formule sans quantificateurs. Alors $\cM \models \phi(\abar) \iff \cN \models \phi(\abar)$
\end{prop}

\begin{proof}
On commence en montrant par induction que l'interpretation des termes dans ces deux structures est identique,
c'est à dire que si $t(\vbar)$ est un terme et $\bbar \in M$, alors $t^{\cM}(\bbar) = t^{\cN}(\bbar)$\\
Si $t$ est une constante $c$, alors  $t^{\cM}(\bbar) = c^{\cM} = c^{\cN} = t^{\cN}(\bbar)$\\
Si $t$ est une variable $v_i$, alors $t^{\cM}(\bbar) = b_i = t^{\cN}(\bbar)$\\
Si $t = f(t_1,\dots,t_n)$ tel que $f$ est une fonction d'arité n, et $t_1,\dots,t_n$ sont des termes tels que pour tout $i \in \{ 1,\dots,n \}$,
$t^{\cM}_i(\bbar) = t^{\cN}_i(\bbar)$, alors
\begin{align*}
t^{\cM}(\bbar) &= f^{\cM}(t^{\cM}_1(\bbar),\dots,t^{\cM}_n(\bbar))\\
&= f^{\cN}(t^{\cM}_1(\bbar),\dots,t^{\cM}_n(\bbar)) \text{ car $f^{\cM} = f^{\cN} \vert_{\cM^n}$,}\\
&= f^{\cN}(t^{\cN}_1(\bbar),\dots,t^{\cN}_n(\bbar)) \text{ par hypothèse de récurrence,}\\
&= t^{\cN}(\bbar) \text{ par définition.}
\end{align*}
On termine la preuve par induction sur les formules:\\
Si $\phi$ est $t_1 = t_2$, alors
\begin{align*}
\cM \models \phi(\abar) &\iff t^{\cM}_1(\abar) = t^{\cM}_2(\abar)\\
&\iff t^{\cN}_1(\abar) = t^{\cN}_2(\abar)\\
&\iff \cN \models \phi(\abar)
\end{align*}
Si $\phi$ est $\cR(t_1,\dots,t_n)$, où $\cR$ est une relation d'arité n, alors
\begin{align*}
\cM \models \phi(\abar) &\iff (t^{\cM}_1(\abar),\dots,t^{\cM}_n(\abar)) \in \cR^{\cM}\\
&\iff (t^{\cM}_1(\abar),\dots,t^{\cM}_n(\abar)) \in \cR^{\cN} \text{ car $\cR^{\cM} = \cR^{\cN} \vert_{\cM^n}$,}\\
&\iff (t^{\cN}_1(\abar),\dots,t^{\cN}_n(\abar)) \in \cR^{\cN}\\
&\iff \cN \models \phi(\abar)
\end{align*}
Si $\phi$ est $\neg\psi$ et qu'on suppose que le résultat est vrai pour $\psi$, alors
\begin{align*}
\cM \models \phi(\abar) &\iff \cM \not\models \psi(\abar)\\
&\iff \cN \not\models \psi(\abar) \text{ par hypothèse}\\
&\iff \cN \models \phi(\abar)
\end{align*}
Si $\phi$ est $\psi_1 \wedge \psi_2$ et qu'on suppose que le résultat est vrai pour $\psi_1$ et $\psi_2$, alors
\begin{align*}
\cM \models \phi(\abar) &\iff \cM \models \psi_1(\abar) \text{ et } \cM \models \psi_2(\abar)\\
&\iff \cN \models \psi_1(\abar) \text{ et } \cN \models \psi_2(\abar) \text{ par hypothèse}\\
&\iff \cN \models \phi(\abar)
\end{align*}

\noindent Le résultat est donc vrai pour toutes les formules atomiques et est stable sous négation et conjonction logique.
Par construction des formules d'un langage (définition \ref{termeformule}), il est vrai pour toutes formules sans quantificateurs.
\end{proof}

\begin{defi}[Structures élémentairement équivalentes]
On dit que deux $\cL$-structures $\cM$ et $\cN$ sont élémentairement équivalentes, et on note $\cM \equiv \cN$, si
pour toute $\cL$-phrases $\phi$, on a
\begin{align*}
\cM \models \phi \iff \cN \models \phi
\end{align*}
\end{defi}

\begin{theo}\label{isomorphisme}
Soient $\mathcal{M}$ et $\mathcal{N}$ deux $\mathcal{L}$-structures. Supposons qu'il existe $j: \mathcal{M} \rightarrow \mathcal{N}$ un isomorphisme,
alors $\mathcal{M} \equiv \mathcal{N}$.
\end{theo}

\begin{proof}
Nous commençons par montrer par induction que les définitions \ref{plongement} et \ref{defi-terme} suffisent à ce que les termes se comportent bien
sous plongement quelconque, i.e pour $t$ un $\mathcal{L}$-terme, muni de variables libres dans $\overline{v} = (v_1,\dots,v_n)$,
et $\overline{a}=(a_1,\dots,a_n) \in M^n$, on a $j(t^{\mathcal{M}}(\overline{a}))
= t^{\mathcal{N}}(j(\overline{a}))$, en notant $j(\overline{a})=(j(a_1),\dots,j(a_n))$\\[0.5cm]
    Si $t=c$ avec $c$ constante, on a 
    \begin{align*}
    j(t^{\cM}(\overline{a})) &= j(c^{\cM})\\
    &= c^{\cN} && \text{par définition d'un plongement,} \\
    &= t^{\cN}(j(\overline{a}))
    \end{align*}
    Si $t=v_i$,\[
    j(t^{\cM}(\overline{a}))  = j(a_i) = t^{\cN}(j(\overline{a}))\]
    Si $t=f(t_1, \ldots, t_r)$, alors
    \begin{align*}
    j(t^{\cM}(\overline{a})) &= j(f^{\cM}(t_1^{\cM}(\overline{a}),\ldots, t_r^{\cM}(\overline{a}))) \\
    & = f^{\cN}(j(t_1^{\cM}(\overline{a}),\ldots,t_r^{\cM}(\overline{a}))) && \text{par définition d'un plongement,} \\
    & = f^{\cN}(t_1^{\cN}(j(\overline{a})),\ldots, t_r^{\cN}(j(\overline{a}))) && \text{par hypothèse de récurrence,} \\
    & = t^{\cN}(j(\overline{a})).
    \end{align*}

Nous terminons la preuve par induction sur les formules:\\[0.5cm]
Si $\phi(\overline{v}): t_1 = t_2$,
\begin{align*}
\mathcal{M} \models \phi(\overline{a}) &\iff t_1^{\mathcal{M}}(\overline{a}) = t_2^{\mathcal{M}}(\overline{a})\\
&\iff j(t_1^{\mathcal{M}}(\overline{a})) = j(t_2^{\mathcal{M}}(\overline{a})) && \text{par injectivité de j}\\
&\iff t_1^{\mathcal{N}}(j(\overline{a})) = t_2^{\mathcal{N}}(j(\overline{a}))\\
&\iff \mathcal{N} \models j(\phi(\overline{a}))
\end{align*}
Si $\phi(\overline{v}): \mathcal{R}(t_1,\dots,t_n)$,
\begin{align*}
\mathcal{M} \models \phi(\overline{a}) &\iff (t_1^{\mathcal{M}}(\overline{a}),\dots,t_n^{\mathcal{M}}(\overline{a})) \in \mathcal{R}^{\mathcal{M}}\\
&\iff (j(t_1^{\mathcal{M}}(\overline{a})),\dots,j(t_n^{\mathcal{M}}(\overline{a}))) \in \mathcal{R}^{\mathcal{N}}\\
&\iff (t_1^{\mathcal{N}}(j(\overline{a})),\dots,t_n^{\mathcal{N}}(j(\overline{a}))) \in \mathcal{R}^{\mathcal{N}}\\
&\iff \mathcal{N} \models \phi(j(\overline{a}))
\end{align*}
Si $\phi: \neg\psi$, pour $\psi$ une certaine $\mathcal{L}$-formule,
\begin{align*}
\mathcal{M} \models \phi(\overline{a}) &\iff \mathcal{M} \not\models \psi(\overline{a})\\
&\iff \mathcal{N} \not\models \psi(j(\overline{a})) && \text{par induction sur } \psi\\
&\iff \mathcal{N} \models \phi(j(\overline{a}))
\end{align*}
Si $\phi: \psi \wedge \theta$,
\begin{align*}
\mathcal{M} \models \phi(\overline{a}) &\iff \mathcal{M} \models \psi(\overline{a}) \text{ et } \mathcal{M} \models \theta(\overline{a})\\
&\iff \mathcal{N} \models \psi(j(\overline{a})) \text{ et } \mathcal{N} \models \theta(j(\overline{a})) && \text{par induction sur }\psi\text{ et }\theta\\
\end{align*}
Si $\phi(\overline{v}): \exists w \text{ } \psi(\overline{v},w)$,
\begin{align*}
\mathcal{M} \models \phi(\overline{a}) &\iff \mathcal{M} \models \psi(\overline{a},b) \text{, pour un certain b dans M,}\\
&\iff \mathcal{N} \models \psi(j(\overline{a}),c) \text{, pour un certain c dans N,  car j est surjective,}\\
&\iff \mathcal{N} \models \phi(j(\overline{a}))
\end{align*}
\end{proof}

\subsection{Théories}

\begin{defi}[Théories, modèles, satisfiabilité]
Soit $\cL$ un langage.
\begin{enumerate}
\item Une $\cL$\kw{-théorie} $T$ est un ensemble de $\cL$-phrases (i.e un ensemble de $\cL$-formules sans variables libres), qu'on appelle \kw{axiomes}.
\item On dit que $\cM$ est un \kw{modèle} de $T$, et on note $T \models \cM$, si l'on a $\cM \models \phi$ pour tout axiome $\phi \in T$.
\item Une théorie $T$ est dite \kw{satisfiable} s'il existe une $\cL$-structure $\cM$ tel que $\cM$ est un modèle de $T$.
\end{enumerate}
\end{defi}

\begin{exem}
On axiomatise la structure de groupe via le langage $\cL = \{ \star , e \}$ et la théorie suivante:
\begin{enumerate}
\item $\forall x,y,z, (x \star y) \star z = x \star (y \star z)$
\item $\forall x, x \star e = e \star x = x$
\item $\forall x, x \star y = y \star x = e$
\end{enumerate}
Cette théorie est satisfiable: par exemple, $\mathbb{R}$ muni de la loi d'addition $+$ et du neutre $e_{\mathbb{R}} = 0$ en est un modèle.
\end{exem}

\begin{exem}
L'objet d'étude central du mémoire est la \kw{théorie des corps algébriquement clos},
dont nous présentons l'axiomatisation dans la définition \ref{defACF}.
\end{exem}

\section{Ensembles définissables}
\begin{defi}[Ensembles définissables]
Soit $\mathcal{M} = (M,\dots)$ une $\mathcal{L}$-structure. On dit qu'un sous-ensemble $X \subseteq M^n$ est \kw{définissable} s'il existe une 
$\mathcal{L}$-formule $\phi(v_1,\dots,v_n,w_1,\dots,w_m)$ et $\overline{b} \in M^m$ tels que $X= \{ \overline{a} \in M^n \mid \mathcal{M} \models 
\phi(\overline{a},\overline{b}) \}$. On dit que $\phi(\overline{v},\overline{b})$ \kw{définit} $X$. Soit $A \subseteq 
M$, on dit que $X$ est \kw{$A$-définissable}, ou \kw{définissable sur $A$}, si le paramètre choisi peut être pris dans $A$, c'est-à-dire s'il
existe une formule $\phi(\overline{v},w_1,\dots,w_l)$ et $\overline{b} \in A^l$, tels que $\phi(\overline{v},\overline{b})$ définit X.
\end{defi}

\begin{exem}
On se munit du langage des anneaux $\mathcal{L}_{ring} = \{0, 1, +, \times\}$ et de son modèle $\mathbb{C}$, le corps des nombres complexes.
Pour $n \in \mathbb{N}^*$, on définit les racines n-ième de l'unité $U_n = \{z \in \mathbb{C}, z^n = 1 \}$.
$U_n$ est alors définissable via la formule 
\begin{align*}
\phi(z): \underbrace{z \times \dots \times z}_{n \text{ fois}} = 1
\end{align*}
Nous précisons que $U_n$ est $\kw{0-définissable}$, c'est à dire que la formule $\phi$ qui la définit n'a pas de paramètre.
\end{exem}

La fin de cette section illustre un premier exemple de résultat en Théorie des Modèles applicable à l'Algèbre.

\begin{prop}\label{pointwise}
Soit $\mathcal{M}$ une $\mathcal{L}$-structure. Si $X \subset M^n$ est $A$-définissable, alors tout $\mathcal{L}$-automorphisme $\sigma$
qui fixe point par point $A$ (i.e pour tout $a \in A$, $\sigma(a) = a$), laisse $X$ invariant (i.e $\sigma(X)=X$)
\end{prop}

\begin{proof}
Soit $\overline{a} \in A$ et $\phi(\overline{v}, \overline{a})$ la $\mathcal{L}$-formule définissant $X$.
Soient $\sigma$ un automorphisme de $\mathcal{M}$ tel que $\sigma(\overline{a}) = \overline{a}$ et $\overline{b} \in M^n$.
Nous avons, 
\begin{align*}
\mathcal{M} \models \phi(\overline{b}, \overline{a}) &\iff \mathcal{M} \models \phi(\sigma(\overline{b}), \sigma(\overline{a})) 
&& \text{par Théorème } \ref{isomorphisme}\\
&\iff \mathcal{M} \models \phi(\sigma(\overline{b}), \overline{a})
\end{align*}
On obtient bien $\bbar \in X \iff \sigma(\bbar) \in X$.
\end{proof}

\begin{coro}
Les nombres réels $\mathbb{R}$ en tant que sous-ensemble du corps des complexes $\mathbb{C}$ n'est pas définissable.
\end{coro}

\begin{proof}
Supposons par l'absurde que $\mathbb{R}$ soit définissable sur $A \subset \mathbb{C}$, en particulier on peut prendre $A$ fini.
$A$ étant fini, sa cloture algébrique est dénombrable: par des arguments de cardinalité, on peut choisir $r,\text{ }s \in \mathbb{C}$ 
algébriquement indépendant sur A et tel que $r \in \mathbb{R}$, $s \not\in \mathbb{R}$.
On a alors l'existence d'un automorphisme $\sigma$ de $\mathbb{C}$ tel que pour tout $x \in A, \sigma(x)=x$ et $\sigma(r)=s$,
ainsi $\sigma(\mathbb{R}) \neq \mathbb{R}$ et par contraposée de la proposition \ref{pointwise}, $\mathbb{R}$ n'est pas définissable sur $\mathbb{C}$.
\end{proof}

\begin{rema}
Bien que la topologie singulière de $\mathbb{R}$ nous pousserait à croire qu'il occupe une place particulière dans $\mathbb{C}$, la Théorie des Modèles
illustre ici que le langage pur des corps est "aveugle" façe à nos intuitions sur la droite réelle.
\end{rema}

\section{Ultraproduits et Théorème de Compacité}

\noindent Cette section s'écarte du plan de Marker pour fournir une preuve sémantique du théorème de compacité fondée sur les ultraproduits.
Tandis que l'approche syntaxique découle du théorème de complétude, nous proposons ici de construire un tel modèle satisfaisant $T$ via une famille de structures.
Nous commençons par énoncer le résultat.

\begin{theo}[Théorème de compacité]
Soit $T$ une $\cL$-théorie tel que tout sous-ensemble fini de $T$ est satisfiable, alors $T$ est satisfiable.
\end{theo}

\subsection{Ultrafiltres}

\begin{defi}[Filtres, ultrafiltres]
Soit $I$ un ensemble quelconque, un filtre sur $I$ est un sous-ensemble $\mathcal{F}$ de $\mathcal{P}(I)$ tel que 
\begin{enumerate}
\item $I \in \mathcal{F}, \emptyset \not\in \mathcal{F}$
\item $X, Y \in \mathcal{F} \implies X \cap Y \in \mathcal{F}$
\item $X \in \mathcal{F} \wedge X \subset Y \implies Y \in \mathcal{F}$
\end{enumerate}
Un filtre est un ultrafiltre s'il est un filtre maximal au sens de l'inclusion.
\end{defi}

\begin{defi}
Soient $I$ un ensemble et $A$ un ensemble de parties de $I$. On dit que $A$ a la propriété de l'intersection finie si,
pour tout $A_0,\dots,A_n \in A$, on a $A_0 \cap \dots \cap A_n \ne \emptyset$.
\end{defi}

\begin{prop}
Soit $A \subset I$ non vide et ayant la propriété de l'intersection finie. L'ensemble des parties de $I$ qui contiennent une intersection finie 
d'éléments de $A$ est un filtre, qu'on appelle le filtre engendré par $A$.
\end{prop}

\begin{proof}
Soit un tel sous-ensemble $A \subset I$. On montre que $ \cF = \{ X \subset I \mid \exists A_0,\dots,A_n \in A, A_0 \cap \dots \cap A_n \in X \}$ est un filtre.
Clairement $I \in \cF$ et $\emptyset \not\in \cF$ par propriété d'intersection finie. Si $X \in \cF$ et $X \subset Y$, alors il existe $A_0,\dots,A_n \in A$
tel que $A_0 \cap \dots \cap A_n \subset X$, d'où $A_0 \cap \dots \cap A_n \subset Y$, et $Y \in \cF$. Si $X,Y \in \cF$, alors il existe $I_X$ et $I_Y$ respectivement
des intersections finies de parties de $A$ dans $X$ et $Y$, et $I_X \cap I_Y$ est encore une intersection finie d'éléments de $A$ tel que $I_X \cap I_Y \subset X \cap Y$,
d'où $X \cap Y \in \cF$.
\end{proof}

\begin{prop}\label{lemmelos}
Soit $I$ un ensemble. Un filtre $\cF$ est un ultrafiltre si et seulement si pour tout $X \subset I$, on a $X \in \cF$ ou bien $I \setminus X \in \cF$.
Tout filtre est contenu dans un ultrafiltre.
\end{prop}

\begin{proof}
\begin{enumerate}
\item On commence par remarquer qu'un filtre $\cF$ sur $I$ ne peut pas contenir une partie $A$ et son complémentaire $I \setminus A$,
sinon $A \cap (I \setminus A) = \emptyset \in \cF$. Supposons alors un filtre $\cF$ sur $I$ tel que $A \not\in \cF$, on montre que
$\cF \cup \{A\}$ a la propriété d'intersection finie si et seulement si $I \setminus A \not\in \cF$. Ceci impose l'implication directe: si $\cF$ est
un ultrafiltre, alors on ne peut pas l'agrandir via $A$, et on obtient $I \setminus A \in \cF$.
\begin{enumerate}[label= (\roman*)]
\item Si $I \setminus A \in \cF$, alors $A \cup (I \setminus A) = \emptyset$ et $\cF \cup \{A\}$ n'a pas la propriété d'intersection finie.
\item Si $\cF \cup \{A\}$ n'a pas la propriété d'intersection finie,
alors il existe $B_0,\dots,B_n \in \cF$ tel que $B_0 \cap \dots \cap B_n \cap A = \emptyset$, d'où $B_0 \cap \dots \cap B_n \subset I \setminus A$,
et $I \setminus A \in \cF$ par stabilité par sur-ensemble.
\end{enumerate}
Pour l'implication réciproque, on suppose par l'absurde qu'il existe un filtre $\cF$ sur $I$ vérifiant : pour tout $A \in I$, $A \in \cF$ ou bien
$I \setminus A \in \cF$ mais que $\cF$ n'est pas maximal. Ainsi on a une partie $B \subset I$ tel que $B \not\in \cF$ et $\cF \cup \{B\}$ est toujours un filtre.
Or $I \setminus B \in \cF$ par hypothèse, d'où $I \setminus B \in \cF \cup \{B\}$. Par stabilité sur intersection, $B \cap (I \setminus B) = \emptyset \in \cF \cup \{B\}$, 
c'est absurde.
\item Pour le deuxième point, invoquons le Lemme de Zorn. Soit $\cF$ un filtre sur $I$. Considérons l'ensemble de tous les filtres sur $I$ contenant $\cF$ ordonné par 
la relation d'inclusion. Soit $(\cF_j)_{j \in J}$ une chaîne d'un tel ensemble. Ainsi $\bigcup_{j \in J} \cF_j$ est un majorant de cette chaîne, et les conditions du
Lemme de Zorn sont remplis: un tel ensemble possède un élément maximal.
\end{enumerate}
\end{proof}

\begin{exem}
Un filtre $\cF$ sur $I$ est dit principal s'il existe $X_0 \subset I, X_0 \ne \emptyset$, tel que $\cF = \{ Y \subset I \mid X \subset Y \}$.
Un filtre principal $\cF$ est un ultrafiltre sur $I$ si et seulement s'il est généré par un élément, c'est à dire s'il existe $x_0 \in I$
tel que $\cF = \{ X \subset I \mid \{x_0\} \in X \}$.
\end{exem}

\subsection{Ultraproduits}

\noindent Soient $I$ un ensemble d'indices, $(A_i)_{i \in I}$ une famille d'ensembles non vide et $\mathcal{A} = \prod_{i \in I} A_i$ leur produit cartésien,
et $a,b \in \mathcal{A}$. Soit $\mathcal{F}$ un filtre sur $I$. On définit une relation sur $\cA$ :
\begin{align*}
a \equiv_{\cF} b \text{ ssi } \{ i \in I \mid a(i) = b(i) \} \in \cF
\end{align*}
où $a(i)$ désigne la $i$-ème coordonnées de $a$, c'est à dire $a(i) \in A_i$. Il découle aisément de la définition d'un filtre que $\equiv_{\cF}$
est une relation d'équivalence. On note $a_{\cF}$ la classe de $a$ sur cette relation.
On introduit alors la définition suivante, en gardant les mêmes notations :

\begin{defi}[Ultraproduits]
Le quotient $\cA / \cF = \{ a_{\cF} \mid a \in \cA \}$ est appellé \textbf{produit réduit} de $\cA$. Dans le cas où $\cF$ est un ultrafiltre,
on appelle ce quotient un \textbf{ultraproduit}. Dans ce dernier cas, si $a \equiv_{\cF} b$, on dit que $a$ et $b$ sont égaux presque partout.
\end{defi}

\begin{defiprop}[Structure de l'ultraproduit]
Soit $I$ un ensemble d'indices. Soient $(M_i)_{i \in I}$ une famille de $\cL$-structures et $\cM = \prod_{i \in I} M_i$ leur produit cartésien.
Soit $\mathcal{U}$ un ultrafiltre sur $I$. L'ultraproduit $\cM/\mathcal{U}$ est naurellement muni d'une $\cL$-structure, en intérpretant les symboles de $\cL$ ainsi:
\begin{itemize}
\item Pour chaque symbole de constante $c \in \cL$, $c^{\cM/\mathcal{U}} = (c^{M_i})_{i \in I}/\mathcal{U}$
\item Pour chaque symbole de fonction $f \in \cL$ d'arité $n$, $f^{\cM/\mathcal{U}}(\abar_{\mathcal{U}}) = (f^{M_i}(\abar(i)))_{i \in I}/\mathcal{U}$
\item Pour chaque relation $R \in \cL$, d'arité $n$, $\cM/\mathcal{U} \models R(\abar_{\mathcal{U}}) \iff \{ i \in I \mid
M_i \models R(\abar(i))\} \in \mathcal{U}$
\end{itemize}
\end{defiprop}

\begin{proof}
On doit montrer que nos interprétations sont bien définies, c'est à dire que le résultat ne dépend pas du choix des représentants des classes d'équivalence.
On détaille le cas des symboles de fonctions.
Supposons $\abar \equiv_{\mathcal{U}} \bbar$. On a ainsi que $I_0 = \{ i \in I \mid \abar(i) = \bbar(i) \} \in \mathcal{U}$.
On doit montrer que $f^{\cM/\mathcal{U}}(\abar) \equiv_{\mathcal{U}} f^{\cM/\mathcal{U}}(\bbar)$, c'est à dire,
$I_1 = \{ i \in I \mid f^{M_i}(\abar(i)) = f^{M_i}(\bbar(i)) \} \in \mathcal{U}$. Or par définition d'une application,
on a clairement $I_0 \subset I_1$, et $\mathcal{U}$ est stable par sur-ensemble, d'où $I_1 \in \mathcal{U}$.
\end{proof}


\begin{theo}[Théorème de \L{}oś]
Soient $I$ un ensemble d'indices et $(M_i)_{i \in I}$ une famille de $\cL$-structures. Soit $\mathcal{U}$ un ultrafiltre sur $I$.
Alors pour $\phi(\xbar)$ une $\cL$-formule et $\abar \in \prod_{i \in I} M_i = \mathcal{M}$,\\on a
\begin{align*}
\mathcal{M}/\mathcal{U} \models \phi(\abar_{\mathcal{U}}) \iff \{ i \in I \mid M_i \models \phi(\abar(i)) \} \in \mathcal{U}
\end{align*}
où $\bar{a}_\mathcal{U}$ désigne la classe d'équivalence du tuple $\bar{a}$ dans l'ultraproduit.
\end{theo}

\begin{rema}
Ce résultat pivot, que l'on admettra, repose entre autre sur la caractérisation des ultrafiltres données dans la proposition \ref{lemmelos}.
Il permet une sorte de "passage à l'infini" de la satisfiabilité d'une formule si et seulement si
cette formule est "vraie presque partout" au sens de l'ultrafiltre. Cette caractérisation est la clé du Théorème de compacité.
\end{rema}

\begin{theo}[Théorème de compacité]\label{theo-compacite}
Soit $T$ une $\cL$-théorie tel que tout sous-ensemble fini de $T$ est satisfiable, alors $T$ est satisfiable.
\end{theo}

\begin{proof}
Soit $T$ une telle théorie. Pour chaque sous-ensemble fini $F$ de $T$, on appelle $\cM_F$ un modèle satisfaisant $F$.
Soit $I$ l'ensemble des parties finies de $T$. Pour chaque axiome $\phi$ de $T$, soit $V(\phi) = \{F \in I \mid \phi \in F \}$
l'ensemble des parties finies de $T$ ayant $\phi$ pour axiome. On remarque que la famille $\{V(\phi) \mid \phi \in T\}$ a la propriété 
de l'intersection finie, en effet pour $T_0 = \{\phi_1,\dots,\phi_n\}$ une famille finie d'axiome, on a $T_0 \in V(\phi_1) \cap \dots \cap V(\phi_n)$ 
car $T_0 \in V(\phi_1) \text{ via } \phi_1, T_0 \in V(\phi_2) \text{ via } \phi_2$ etc. Ainsi $\{V(\phi) \mid \phi \in T\}$ engendre un filtre,
donc est contenue dans un ultrafiltre $\mathcal{U}$. On pose $\cM = \prod_{F \in I} \cM_F / \mathcal{U}$ l'ultraproduit des $\cM_F$. On montre que $\cM$ est un modèle de $T$.
Soit $\phi \in T$, on a $\cM \models \phi$ si et seulement si $J_{\phi} = \{ F \in I \mid \cM_F \models \phi \} \in \mathcal{U}$ par Théorème de \L{}oś.
\begin{enumerate}
\item Si $F \in V(\phi)$, alors trivialement $\cM_F \models \phi$, d'où $V(\phi) \subset J_{\phi}$.
\item Or $V(\phi) \in \mathcal{U}$ par construction, et $\mathcal{U}$ est stable par sur-ensemble, d'où $J_\phi \in \mathcal{U}$.
\end{enumerate}
Nous avons construit un modèle de $T$: $\cM \models T$.
\end{proof}

\section{Complétude, categoricité et test de Vaught}

\begin{defi}[Théorie complète]
Une $\mathcal{L}$-théorie $T$ est dite complète si pour toute $\mathcal{L}$-phrase $\phi$, on a soit $T \models \phi$, soit $T \models \neg\phi$.
\end{defi}

\begin{defi}[Categoricité]
Soient $\kappa$ un cardinal infini et $T$ une $\cL$-théorie satisfaite par des modèles de cardinalité $\kappa$.
On dit que $T$ est $\kappa$-catégorique si les modèles de $T$ de cardinalité $\kappa$ sont deux à deux isomorphes.
\end{defi}

\noindent Nous donnons un premier exemple de théorie $\kappa$-catégorique: la théorie au langage vide forçant ses modèles à être infini.

\begin{prop}
Soit $\cL = \emptyset$. Pour $n \in \mathbb{N^*}$, on définit la $\cL$-phrase
\begin{align*}
\phi_n = \exists x_1\dots\exists x_n(\bigwedge_{1\leq i < j \leq n} x_i \neq x_j)
\end{align*}
On définit la $\cL$-théorie $T_{\infty} = \{\phi_n,\text{ }n \in \mathbb{N^*}\}$.\\
Alors pour tout cardinal $\kappa$, $T_{\infty}$ est $\kappa$-catégorique.
\end{prop}

\begin{proof}
Soit $\kappa$ un cardinal infini. Soient $\cM$ et $\cN$ deux modèles de $T_{\infty}$ tel que $\lvert \cM \rvert = \lvert \cN \rvert = \kappa$.
Puisque $\cL = \emptyset$, nos modèles sont simplement des ensembles. Or par un argument ensembliste, $M$ et $N$ étant de même cardinalité, il existe $\phi: M \longrightarrow N$ une bijection.
$\phi$ étant une bijection entre structures vides, il n'y a aucun symbole à preserver et c'est un $\cL$-isomorphisme.
\end{proof}

\begin{rema}
Ce premier résultat permet de nous familiariser avec la notion de categoricité, mais aussi de comprendre que c'est bien l'ajout d'une structure
algébrique qui permet de distinguer des modèles de même cardinalité via des invariants. Nous étudions maintenant le cas de la théorie $\ACF_p$, qui perd la catégoricité pour $\kappa = \aleph_0$.
\end{rema}


\begin{prop}
Soit $p \in \{0\} \cup \mathcal{P}$, alors $ACF_p$ est $\kappa$-catégorique pour tout cardinal $\kappa$ indénombrable.
\end{prop}

\begin{prop}
Soit $T$ une $\cL$-théorie satisfaite par des modèles infinis. Si $\kappa$ est un cardinal infini tel que $\kappa > \lvert \cL \rvert$, alors
il existe un modèle de $T$ de cardinalité $\kappa$.
\end{prop}

\begin{theo}[Test de Vaught]
Soit $T$ une $\cL$-théorie satisfaisable, sans modèles finis, et $\kappa$-catégorique pour un certain cardinal infini $\kappa > \lvert \cL \rvert$.
Alors $T$ est complète.
\end{theo}

\begin{coro}\label{coro-acf-complete}
Soit $p \in \{0\} \cup \mathcal{P}$, alors $ACF_p$ est complète.
\end{coro}

\todo[inline]{sous section sur degré de transcandance et preuves}
\todo[inline]{appeler steinitz, donner exemples pour demystifier ces manipulations d'infini}
\todo[inline]{réfléchir sur ce qu'on met de ACF dans cette section et ce qu'on met dans chap2}

\todo[inline]{}\label{lemm-sous-struct-sans-quant}
\todo[inline]{}\label{prop-sous-struct-existance}
\todo[inline]{}\label{defi-diagramme}
\todo[inline]{}\label{defi-decidable}
